\chapter{Literature Survey}

IoT authentication research has increasingly favoured lightweight and distributed
designs, driven by the prohibitive load that centralised systems impose on a single
gateway and the single point of failure they introduce.
Given the severe power and memory constraints of IoT devices, the research community
has sought architectures in which multiple network nodes share the authentication
workload.

\textbf{Light-SAE}~\cite{rosa2023lightsae} exemplifies this trend.
It enables already-authenticated nodes to assist newcomers through a simple
challenge-and-response join protocol, using TinyJAMBU for lightweight encryption
and a modified Diffie--Hellman exchange to establish the session key in multihop
networks, thereby reducing the load on the gateway.
However, Light-SAE transmits node identities in plaintext and provides no
protection for anonymity or session unlinkability, rendering it unsuitable for
privacy-sensitive deployments.

\textbf{SAKMS}~\cite{yang2024sakms} targets industrial IoT systems operating under
the 6TiSCH schedule, employing ECC-based mutual authentication with ECQV
certificates, formalised key management, and strong forward secrecy, with
correctness confirmed by ProVerif.
Its reliance on multiple elliptic-curve operations, however, imposes overhead that
exceeds what severely constrained IoT nodes can sustain.
SAKMS also embeds fixed device identifiers within messages, leaving devices
susceptible to tracking across sessions.
These limitations illustrate that cryptographic strength alone does not imply
privacy preservation or lightweight execution.

The \textbf{DAuth PUF-based authentication scheme}~\cite{das2026comsnets}
addresses many of these weaknesses.
It eliminates long-term secret storage by deriving device secrets from PUFs, uses
a hash-based accumulator for membership verification, and decouples the
authentication server from the parent node so that both parent and non-parent nodes
can serve as authenticators, distributing load in multihop networks.
Its computational primitives (hash, XOR, and two AES calls) keep the
per-device footprint minimal.
Its primary limitation is that the device identity ($ID_D$) is transmitted directly
in every authentication exchange; since all messages traverse a public channel, a
passive observer can capture the identity, track the device over time, and mount
targeted attacks.
Identity concealment is therefore a necessary extension to this protocol.

Several anonymous authentication schemes have been proposed to address such privacy
concerns.
\textbf{He~et~al.}~\cite{he2023lightweight} achieve lightweight anonymous
authentication using hash-and-XOR primitives for CoAP/OSCORE but rely on a
centralised server, forfeiting the decoupled-distribution advantage.
\textbf{LAAKA}~\cite{ali2024laaka} protects device identity during mutual
authentication using only lightweight hash and XOR operations; its overhead,
however, stems from larger message payloads and additional message exchanges
on severely constrained nodes.

A comprehensive survey on lightweight anonymous
authentication~\cite{zhong2025survey} catalogues approaches including temporary
pseudonyms, PUF-assisted anonymous login, hidden membership proofs via accumulators,
and unlinkable tokens, while documenting common pitfalls such as desynchronisation,
pseudonym exhaustion, and inadequate replay protection.
Their taxonomy deconstructs existing protocols into nine sub-frameworks (SF1-UA
through SF9-FS) and two general frameworks (GF1, GF2), finding that only GF2
simultaneously achieves user anonymity, forward secrecy, and desynchronisation
resistance.

Building on this taxonomy, the present work adopts \textbf{GF2} as its target
architectural pattern and realises a \emph{modified} instantiation of it tailored to
resource-constrained multihop IoT nodes.
Concretely, the proposed scheme combines the user-anonymity sub-frameworks
(SF3-UA, SF5-UA, SF6-UA), which use rotating pseudonyms instead of a static
identity, with the forward-secrecy sub-frameworks (SF7-FS, SF8-FS, SF9-FS), which
derive each session key from independently refreshed material.
Unlike the survey's reference constructions that rely on heavy cryptographic
primitives, the proposed scheme uses lightweight PUF-based secret derivation and
hash-XOR masking, and meets GF2's desynchronisation-resistance requirement through
a dual-state recovery mechanism at the authentication server.
The result is a GF2-compliant design with the same low overhead as
DAuth~\cite{das2026comsnets}.

\textbf{Zhou~et~al.}~\cite{zhou2024auth} propose a security-enhanced lightweight
and anonymity-preserving user authentication scheme for IoT-based healthcare,
providing hash-based anonymous authentication with nonce masking to protect sensitive
metadata over open wireless channels.

Other works use \textbf{blockchain}~\cite{ahmed2023chains} and
\textbf{certificateless cryptography}~\cite{alriyami2003cert} to achieve
decentralisation and identity hiding, but they incur prohibitive storage and consensus
overhead.
\textbf{Li~et~al.}~\cite{li2020puf} and \textbf{Zheng~et~al.}~\cite{zheng2022puf}
combine PUF with session key exchange but rely on ECC operations too expensive for
highly constrained nodes.

Despite this body of work, no existing scheme simultaneously satisfies all of the
following requirements: (i)~decoupled authentication, (ii)~PUF-based secrecy
without long-term key storage, (iii)~low-overhead hash-XOR computation, and
(iv)~device anonymity with session unlinkability and desynchronisation recovery.
This gap motivates the present work, which builds an anonymous and
desynchronisation-resilient extension of DAuth~\cite{das2026comsnets} while
preserving its lightweight and distributed properties.

To show the differences between all these methods clearly, Table~\ref{tab:features}
lists the use of decoupling, PUF, accumulator-based authentication, authentication
secret update, anonymity, desynchronisation recovery, and the computation and
communication overhead for all the schemes studied.

\begin{table}[H]
  \centering
  \caption{Features Comparison of Authentication Schemes.}
  \label{tab:features}
  \small
  \setlength{\tabcolsep}{5pt}
  \begin{tabular}{lcccccccc}
    \toprule
    Scheme & Dec. & PUF & Acc. & ASU & Anon. & Desync & C.OH & M.OH \\
    \midrule
    RAAF-MEC~\cite{alruwaili2025raaf}   & $\times$ & \checkmark & $\times$ & \checkmark & $\times$ & $\times$ & H & L \\
    SAKMS~\cite{yang2024sakms}          & \checkmark & $\times$ & $\times$ & $\times$ & $\times$ & $\times$ & H & L \\
    Li et al.~\cite{li2020puf}          & \checkmark & \checkmark & $\times$ & \checkmark & $\times$ & $\times$ & H & L \\
    Zheng et al.~\cite{zheng2022puf}    & \checkmark & \checkmark & $\times$ & \checkmark & $\times$ & $\times$ & M & L \\
    LAAKA~\cite{ali2024laaka}           & $\times$ & $\times$ & $\times$ & \checkmark & \checkmark & $\times$ & L & L \\
    Zhou et al.~\cite{zhou2024auth}     & $\times$ & \checkmark & $\times$ & \checkmark & \checkmark & $\times$ & H & M \\
    DAuth~\cite{das2026comsnets}         & \checkmark & \checkmark & \checkmark & \checkmark & $\times$ & $\times$ & L & L \\
    \textbf{Proposed}                  & \checkmark & \checkmark & \checkmark & \checkmark & \checkmark & \checkmark & \textbf{L} & \textbf{L} \\
    \bottomrule
  \end{tabular}\\[5pt]
  \raggedright\small
  \textit{Dec.}: Decoupled design.\quad
  \textit{Acc.}: Hash-based accumulator.\quad
  \textit{ASU}: Auth.\ secret update.\\
  \textit{Anon.}: Device anonymity via pseudonym.\quad
  \textit{Desync}: Desync.\ recovery.\\
  \textit{C.OH}: Computation: L (Hash/XOR/$\le$2 AES), M ($>$2 AES), H (ECC/Pairings).\\
  \textit{M.OH}: Message exchanges: L (1--3), M (4--5), H ($>$5).
\end{table}
